\documentclass[12pt,oneside]{memoir}
\usepackage[latinica]{matfmaster}
\usepackage{lmodern}
\usepackage{opensans}
\renewcommand{\headersans}{\opensans}
\usepackage{amsmath,amssymb,mathtools}
\usepackage{booktabs}
\usepackage{enumitem}

\autor{[Ime i prezime studenta]}
\naslov{Generalizovane kvadrike u nastavi i vizuelizaciji matematike}
\godina{[godina odbrane]}
\mentor{[Ime i prezime mentora]}
\komisijaA{[Ime i prezime člana komisije]}
\komisijaB{[Ime i prezime člana komisije]}
\apstr{[Apstrakt će biti dopunjen nakon izrade rada.]}
\kljucnereci{generalizovane kvadrike, geometrijska vizuelizacija, GeoGebra 3D, metodika nastave matematike}
\bib{references}

\begin{document}
\frontmatter
\naslovna
\komisija
\apstrakt
\tableofcontents*

\mainmatter
\chapter{Uvod, motivacija i ciljevi rada}

\chapter{Klasične kvadrike i konusni preseci}

\chapter{Teorijska analiza generalizovanih kvadrika}
\section{Milerova kubična površ i aksijalna inverzija}
\section{Superkvadrike i Lameove krive}
\section{Lokalna teorija površi, Dupinove indikatrise i Plikerov konoid}

\chapter{Metodičko-pedagoški pristup obradi generalizovanih kvadrika}
\section{Didaktička transpozicija i razvoj prostorne intuicije}
\section{Dizajn interaktivnih GeoGebra 3D modela i radionica}
\section{Protokol planirane pedagoške studije}

\chapter{Zaključak i smernice za budući rad}

\appendix
\chapter{Predlozi radnih listova i uputstva za softver}

\backmatter
\literatura
\begin{biografija}
[Biografija autora biće uneta kada budu dostavljeni odgovarajući podaci.]
\end{biografija}
\end{document}
